\documentclass{beamer}
\usetheme{Goettingen}
\usepackage 			[oldstyle,proportional]{ETbb}	% ET Book, oldstyle and proportional
\usepackage{fontawesome}

\title{\LaTeX\ for Primary Source Handouts}
\author{Corey Stephan, Ph.D.}
\institute{\faHome\space{\href{https://www.coreystephan.com}{coreystephan.com}}}

% \texttt{Sample Code}

\begin{document}
\begin{frame}[plain]
    \maketitle
\end{frame}

\begin{frame}
    \frametitle[Everything]{\LaTeX\ for Everything}
The world's best established software system for typesetting and document formatting, \LaTeX\ has been popular throughout the mathematical and physical sciences for decades. The utility of \LaTeX\ for scholars in humanities disciplines, however, tends to be overlooked. \break

\LaTeX\ can be used for typesetting in (almost) all instances in which a scholar of humanities needs to present our materials elegantly. \break

I have prepared these slides in \LaTeX\ -- more quickly and with less frustration than would have been possible in graphical presentation program (e.g. LibreOffice Present).
\end{frame}

\begin{frame}
	\frametitle[Tutorial]{This Tutorial}
In this tutorial, I will introduce attendees to the use of LaTeX for a common task in medievalist scholarship, namely, typesetting texts with unusual line length and other organizational needs, such as primary texts with our original translations. \break

The specific example will be preparing a handout for a historical conference or class presentation with perfectly organized excerpts of a Byzantine Greek text and an original English translation in parallel. \break
\end{frame}

\begin{frame}
	\frametitle[Tutorial]{This Tutorial}
I will walk attendees in sequence through setup, formatting, and, finally, viewing a document of this kind -- talking about different approaches that one might take depending on needs, such as placing right-to-left (notably, Semitic) and left-to-right languages side-by-side. \break

I aim to inspire attendees to think about \LaTeX\ as a highly flexible tool with a wider range of professional applications than it often is understood as having. 
\end{frame}

\end{document}
